
\graphicspath{{./chap2/images/}}
\chapter{Linux}
  
\textbf{이 장은 Linux 명령어, bash script, .bashrc 설정, WSL과 VSCode 연결 및 설정 등에 대해 다룬다.}

\section{Linux의 파일 시스템과 파일 관리}
\subsection{Linux의 파일 시스템 구조}
컴퓨터를 사용하기 위해서 가장 먼저 알아야 할 것은 무엇일까? 여러 가지가 있겠지만 파일이 어디에 있는지 알고, 그 위치로 이동해서 파일을 연 다음에 수정하고 저장하거나, 복사/이동/삭제 등을 할 줄 아는 것이 중요할 것이다. Windows와 같은 GUI 환경에서는 파일 탐색기를 통해 이런 작업을 할 수 있지만 CLI 환경에서는 명령어만으로 이러한 작업을 진행해야 하기 때문에 다소 어려움이 있을 수 있다.

일단 먼저 Linux의 파일 시스템 구조에 대해 알아보자. Windows의 경우, 디스크가 하나만 있다고 가정하면 모든 파일은 \code{C:\} 드라이브 아래에 존재한다. 하지만 Linux에서는 모든 파일이 \code{/}라는 루트 디렉토리 아래에 존재한다. 그리고 그 아래에 여러 디렉토리가 존재하는 트리 구조로 되어 있다. 예를 들어, \code{/home/username/dir1/dir2/file.txt}라는 경로가 있다면, \code{/}는 루트 디렉토리, \code{home}은 루트 디렉토리 아래의 디렉토리, \code{username}은 \code{home} 디렉토리 아래의 디렉토리, \code{dir1}은 \code{username} 디렉토리 아래의 디렉토리, \code{dir2}는 \code{dir1} 디렉토리 아래의 디렉토리, 그리고 \code{file.txt}는 \code{dir2} 디렉토리 아래의 파일이다.

루트 디렉토리 아래에는 여러 다른 디렉토리들이 존재하는데, Linux 시스템의 기본적인 필수 파일들이 있는 디렉토리들도 있기 때문에 조심해야 한다. 이 문서의 수준에서는 보통 사용자가 파일을 저장하는 디렉토리인 \code{/home/username} 아래에서 작업을 진행하게 될 것이다.

리눅스에서 파일이나 디렉토리를 지정할 때는 절대경로와 상대경로를 사용할 수 있다. 절대경로는 \code{/home/username/dir1/dir2}와 같이 루트 디렉토리에서 시작해서 내가 원하는 파일이나 디렉토리의 절대적인 위치이다. 반변, 상대경로는 현재 내 위치에서 시작하는 경로다. 즉, 내가 지금 \code{/home/username/dir1}에 있다면 여기서 \code{/home/username/dir1/dir2}의 상대경로는 그냥 \code{dir2}다.

\subsection{디렉토리 및 탐색}
이제 파일 시스템 구조를 알았으니, 본격적으로 이를 탐색하는 방법에 대해 알아보자. Linux에서는 \code{cd} 명령어를 사용해서 디렉토리를 탐색할 수 있다. \\ \code{.}은 현재 디렉토리, \code{..}은 상위 디렉토리, \code{-}은 이전 디렉토리, \code{~}는 홈 디렉토리를 의미한다. \code{cd <dir>}라 입력하면 해당 디렉토리로 이동된다. 디렉토리를 표현할 때 절대경로, 상대경로 둘 다 가능하며, 실행시 주소쪽 값이 변경된다.
\begin{lstlisting}
    userid:~$ cd dir1
    userid:~/dir1$ cd /home/username/dir1 
    userid:/home/username/dir1$ cd ..
    userid:/home/username$ cd -
    userid:/home/username/dir1$ cd ~
    userid:~$
\end{lstlisting}
또한 \code{pwd} 명령어를 사용해서 현재 위치한 디렉토리를 확인할 수 있다. 

다음으로 배울 명령어는 \code{ls} 명령어이다. 이 명령어는 현재 디렉토리에 있는 파일과 디렉토리 목록을 보여주는 명령어다. 만약 \code{~} 하위에 \code{dir1}이란 디렉토리랑 \code{file1.txt}라는 파일이 있다면, \code{ls} 명령어를 입력했을 때 다음과 같이 출력될 것이다.
\begin{lstlisting}
    userid:~$ ls
    dir1  file1.txt
\end{lstlisting}

\code{ls -l} 명령어를 사용하면 파일과 디렉토리의 상세 정보를 볼 수 있다. 
\begin{lstlisting}
    userid:~$ ls -l
    drwxr-xr-x  2 username username 4096 Mar  6 14:26 dir1
    -rw-r--r--  1 username username    0 Mar  6 14:26 file1.txt
\end{lstlisting}
각각의 열은 파일 또는 디렉토리의 권한, 링크 수, 소유자, 그룹, 크기, 마지막 수정 시간, 그리고 이름을 나타낸다. 권한이나 링크의 경우, 이 문서 수준에서 건드릴 일이 없고, 소유자, 파일 크기, 마지막 수정 시간이 중요할 것이다. 파일 크기는 바이트 단위로 표시된다.

\code{ls -a} 명령어를 사용하면 숨김 파일과 디렉토리도 볼 수 있다. 숨김 파일은 파일 명이 \code{.}으로 시작한다. 마지막으로 \code{ls -la} 명령어를 사용하면 상세 정보와 숨김 파일 모두 볼 수 있다.

\subsection{파일 복사/이동/삭제}

파일을 복사하려면 \code{cp} 명령어를 사용한다. 예를 들어, \code{file1.txt}라는 파일을 \code{file2.txt}로 복사하려면 다음과 같이 입력한다.
\begin{lstlisting}
    userid:~$ cp file1.txt file2.txt
\end{lstlisting}
이렇게 하면 현재 디렉토리에 있는 \code{file1.txt}라는 파일이 \code{file2.txt}라는 이름으로 복사된다. 만약 \code{file2.txt}라는 파일이 이미 존재한다면, \code{cp} 명령어는 기존 파일을 덮어쓰게 된다. 만약 덮어쓰기를 방지하려면 \code{cp -i} 명령어를 사용하면 된다. 이 명령어는 덮어쓰기 전에 사용자에게 확인을 요청한다. 파일을 다른 디렉토리에 복사하고 싶으면 절대경로나 상대경로로 대상 디렉토리를 지정하면 된다. 예를 들어, \code{file1.txt}를 \code{dir1} 디렉토리로 복사하려면 다음과 같이 입력한다.
\begin{lstlisting}
    userid:~$ cp file1.txt dir1/file2.txt
\end{lstlisting}
이렇게 하면 \code{file1.txt}라는 파일이 \code{dir1} 디렉토리에 \code{file2.txt}라는 이름으로 복사된다. \code{cp} 명령어는 디렉토리를 복사할 때는 \code{-r} 옵션을 사용해야 한다. 예를 들어, \code{dir1}이라는 디렉토리를 \code{dir2}라는 이름으로 복사하려면 다음과 같이 입력한다.
\begin{lstlisting}
    userid:~$ cp -r dir1 dir2
\end{lstlisting}

파일을 이동하려면 \code{mv} 명령어를 사용한다. 예를 들어, \code{file1.txt}라는 파일을 \code{file2.txt}로 이동하려면 다음과 같이 입력한다.
\begin{lstlisting}
    userid:~$ mv file1.txt file2.txt
\end{lstlisting}
이렇게 하면 \code{file1.txt}라는 파일이 \code{file2.txt}라는 이름으로 이동된다. \code{mv} 명령어도 \code{cp} 명령어와 비슷하게 덮어쓰기를 하며, 이를 방지할 수도 있고, 다른 디렉토리로 이동시킬 수도 있다. \code{mv} 명령어 역시 디렉토리를 이동시킬 수도 있으며, \code{cp}와 달리 이때는 \code{-r} 옵션이 필요하지 않다.

마지막으로 파일을 삭제하려면 \code{rm} 명령어를 사용한다. 예를 들어, \code{file1.txt}라는 파일을 삭제하려면 다음과 같이 입력한다.
\begin{lstlisting}    
    userid:~$ rm file1.txt
\end{lstlisting}
이렇게 하면 현재 디렉토리에서 \code{file1.txt}라는 파일이 삭제된다. 만약 디렉토리를 삭제하려면 \code{rm -r} 명령어를 사용해야 한다. 예를 들어, \code{dir1}이라는 디렉토리를 삭제하려면 다음과 같이 입력한다.
\begin{lstlisting}
    userid:~$ rm -r dir1
\end{lstlisting}
\code{rm} 명령어는 삭제하기 전에 사용자에게 확인을 요청하는 \code{-i} 옵션도 제공한다. 
\textbf{주의: Linux에는 Windows와 달리 휴지통 개념이 없다. 즉, 파일이나 디렉토리를 삭제한 이후에는 복구하는 것이 쉽지 않다. 따라서 중요한 파일이나 디렉토리를 삭제할 때는 항상 주의해야 한다.}

\subsection{파일 만들기 및 열기}
파일을 만들려면 \code{touch} 명령어를 사용한다. 예를 들어, \code{file1.txt}라는 파일을 만들려면 다음과 같이 입력한다.
\begin{lstlisting}
    userid:~$ touch file1.txt
\end{lstlisting}
이렇게 하면 현재 디렉토리에 \code{file1.txt}라는 이름의 빈 파일이 만들어진다.

파일을 열려면 여러 가지 방법이 있다. 가장 간단한 방법은 \code{cat} 명령어를 사용하는 것이다. 예를 들어, \code{file1.txt}라는 파일의 내용을 보려면 다음과 같이 입력한다.
\begin{lstlisting}
    userid:~$ cat file1.txt
\end{lstlisting}
이렇게 하면 \code{file1.txt}라는 파일의 내용이 터미널에 출력된다. 원래 \code{cat} 명령어는 con\textbf{cat}enate에서 온 명령어로, 여러 파일을 붙이는 역할을 하는 명령어지만, 위와 같이 파일의 내용을 보기 위해서도 사용할 수 있다.

만약 파일이 너무 길어서 한 화면에 다 보이지 않는다면, \code{less} 명령어를 사용할 수 있다. 예를 들어, \code{file1.txt}라는 파일의 내용을 페이지 단위로 보려면 다음과 같이 입력한다.
\begin{lstlisting}
    userid:~$ less file1.txt
\end{lstlisting}
이렇게 하면 \code{file1.txt}라는 파일의 내용이 페이지 단위로 출력된다. \code{less} 명령어를 사용할 때는 엔터를 눌러서 다음 줄로, 스페이스바를 눌러서 다음 페이지로 넘어갈 수 있다. 또한 방향키로도 움직일 수 있다.